\documentclass[11pt,a4paper]{article}
\usepackage{amsmath,amssymb,amsthm}
\usepackage[margin=2.5cm]{geometry}
\usepackage{hyperref}

\newtheorem{definition}{Definition}[section]
\newtheorem{theorem}[definition]{Theorem}
\newtheorem{proposition}[definition]{Proposition}
\newtheorem{lemma}[definition]{Lemma}
\newtheorem{corollary}[definition]{Corollary}
\newtheorem{conjecture}[definition]{Conjecture}
\newtheorem{remark}[definition]{Remark}

\title{Hyperseed Formalization:\\
  The Orchard Bug and the Unfolding Software-Verification Reckoning}
\author{ProtomegaTron (automated formalization)}
\date{2026-09-17}

\begin{document}
\maketitle

\begin{abstract}
We formalize the intellectual structure of Ben Goertzel's Substack article
``The Orchard Bug and the Unfolding Software-Verification Reckoning''
(2026-06-05) within the Hyperseed ontology. The article uses the Zcash
Orchard pool bug as a case study for the coming software-verification
reckoning: AI tools are making bug discovery cheap, forcing the industry
toward formal verification. Each structural insight maps onto Hyperseed
structures: composition failure as cocycle defect, correct-by-construction
as fiber-base identity, AI-accelerated discovery as base space equalization
of vulnerability, formal verification as fiber inspection (proof obligation
= cocycle condition), and the industry's speed preference as provenance debt.
\end{abstract}

\tableofcontents

%% =================================================================
\section{Composition failure as cocycle defect}
\label{sec:cocycle}
%% =================================================================

\begin{definition}[The Orchard bug --- claims 1--3, 20]
\label{def:orchard}
The Orchard pool bug was an under-constrained element in a zero-knowledge
circuit. The key structural feature: individual modules tested fine; the
composition failed because an invariant that held within each module was
not enforced at the interface between them. The bug hid for approximately
four years.

In Hyperseed terms, this is a \emph{cocycle defect}: the transition
functions between modules (the constraints that should hold at module
boundaries) don't compose correctly:
\[
  g_{ij} \circ g_{jk} \neq g_{ik}
\]

The invariant holds locally in each module (each $g_{ij}$ individually
satisfies its local constraint), but the cocycle condition fails globally
--- the composition of boundary constraints doesn't close. The bug
exists precisely in the gap between local correctness and global
consistency.

This is not a random implementation error but a \emph{structural failure}
--- the kind of failure that testing (which checks individual paths)
cannot reliably catch, because the failure only manifests when modules
are composed in ways that testing may not exercise. The bug is invisible
to any analysis that examines modules individually; it lives in the
\emph{transition functions between them}.

The four-year hiding period illustrates a general principle: cocycle
defects are the hardest bugs to find because they require understanding
the global composition structure, not just the local module behavior.
Human auditors tend to verify modules individually (checking local
constraints) rather than verifying the cocycle condition (checking that
constraints compose correctly across all interfaces).
\end{definition}

%% =================================================================
\section{AI-accelerated discovery as base space equalization}
\label{sec:ai-discovery}
%% =================================================================

\begin{proposition}[Economics flip --- claims 4--6, 22]
\label{prop:economics}
AI tools --- code assistants, fuzzing agents, automated auditors ---
are dramatically lowering the cost of finding exploitable bugs. The
article argues this flips the economics: previously, many bugs existed
but finding them cost more than the expected payoff. AI drops the
discovery cost while the deployed codebase stays the same.

In Hyperseed terms, this is \emph{base space equalization of
vulnerability}: AI tools equalize the information base space ---
bugs that were hidden (asymmetric information between defenders and
attackers) become visible (symmetric information):
\[
  \text{AI discovery}: I_{\text{defender}} \approx I_{\text{attacker}}
  \to I_{\text{symmetric}}
\]

This is the same base-space-equalization dynamic as capability
diffusion in the AI-models context (where model capability equalizes
across providers over months), but applied to the vulnerability
landscape (where hidden bugs become discoverable at scale).

The article emphasizes that crypto is just the canary: crypto codebases
are ``newer, smaller, open-source, and written by people who are at
least nominally paranoid'' --- probably in better shape than legacy
banking systems that ``nobody alive fully understands anymore.'' The
vulnerability equalization is coming for \emph{all} software.

The structural consequence: the software industry's accumulated
\emph{provenance debt} (code whose correctness provenance is ``we
tested it'' rather than ``we proved it'') is being called in by
AI-accelerated discovery. The debt was sustainable when discovery
was expensive; it becomes unsustainable when discovery is cheap.
\end{proposition}

%% =================================================================
\section{Formal verification as fiber inspection}
\label{sec:verification}
%% =================================================================

\begin{theorem}[Proof obligation = cocycle condition --- claims 7--12, 23]
\label{thm:verification}
The article argues that formal verification --- proving with mathematics
that software does what it's designed to do --- is the cure, and that
it's well understood, not exotic, and not waiting on future research.

In Hyperseed terms, formal verification is \emph{fiber inspection}:
verifying that the transition functions (module boundary constraints)
compose correctly across all interfaces. The proof obligation IS the
cocycle condition:
\[
  \text{Proof obligation}: \forall i,j,k: \;
  g_{ij} \circ g_{jk} = g_{ik}
\]

If the proof obligation discharges, the cocycle closes --- the
composition of constraints is globally consistent. If it fails,
there's a cocycle defect --- the Orchard bug scenario.

The article identifies three levels of verification practice:
\begin{enumerate}
\item \textbf{Test-and-hope (current norm).} Check individual paths
  (local constraints) without verifying global composition. The
  cocycle condition is \emph{assumed} rather than \emph{verified}.
\item \textbf{Write-then-verify.} Build the implementation, then
  prove it correct. The cocycle condition is checked after
  construction --- a gap exists during development that might
  persist if the proof is never completed.
\item \textbf{Correct by construction.} Derive the implementation
  from the specification. The cocycle condition is satisfied by
  construction --- there is no gap.
\end{enumerate}

The industry has consistently chosen level 1 because levels 2 and 3
cost more effort. AI is forcing the transition to levels 2 and 3 by
making the cost of \emph{not} verifying (bug discovery and exploitation)
higher than the cost of verifying.

The Cardano precedent (formally verifying Ouroboros consensus) shows
that level 2 is practical for real systems. ASI:Chain's approach shows
that level 3 is achievable for critical subsystems.
\end{theorem}

%% =================================================================
\section{Correct by construction as fiber-base identity}
\label{sec:correct}
%% =================================================================

\begin{definition}[Implementation = specification --- claims 13--16, 21, 24]
\label{def:correct}
The article describes ASI:Chain's approach: for key parts of the system,
implementation is derived directly from mathematical specification.
``The implementation and its proof of correctness are two views of the
same underlying mathematical object.''

In Hyperseed terms, this is \emph{fiber-base identity}: the fiber
(intended behavior, specification) and the base (actual behavior,
implementation) are the same mathematical object:
\[
  F_{\text{spec}} \cong B_{\text{impl}}
\]

There is no gap for a cocycle defect to hide in because there is no
separate composition step --- the fiber IS the base. The derivation
process ensures that every property of the specification is
automatically a property of the implementation, and vice versa.

This maps onto the broader Hyperseed distinction between section-level
and fiber-level approaches:
\begin{itemize}
\item \textbf{Verify-after-the-fact = section-level.} You build a
  section (implementation) and then check whether it's consistent
  with the fiber (specification). The section exists independently
  of the fiber, and the verification is a separate step that checks
  their compatibility.
\item \textbf{Correct-by-construction = fiber-level.} You derive the
  section from the fiber, guaranteeing consistency by construction.
  The section is not an independent artifact but a projection of
  the fiber into the implementation space.
\end{itemize}

The concrete example: a variant of the Cordial Miners protocol ---
reputation-based, ledger-free, suited to mesh networks and weak
heterogeneous hardware --- derived correct-by-construction from its
mathematical specification using MeTTa and Rholang (via the
F1R3FLY.io collaboration). The mathematical foundations of these
languages make the derivation tractable.
\end{definition}

%% =================================================================
\section{AI-assisted verification}
\label{sec:ai-verify}
%% =================================================================

\begin{proposition}[Offense-defense symmetry --- claims 17--19]
\label{prop:ai-verify}
The article notes an offense-defense symmetry in AI's relationship
to software correctness: the same AI capabilities that find bugs can
help prove their absence.

Neural-symbolic systems (Hyperon) are particularly well-suited for
AI-assisted formal verification because they combine:
\begin{itemize}
\item \textbf{Neural pattern recognition:} suggesting proof strategies,
  identifying likely invariants, recognizing structural similarities
  to previously verified code.
\item \textbf{Symbolic proof construction:} constructing formal proofs,
  checking proof obligations, verifying that the cocycle condition
  holds across all interfaces.
\end{itemize}

This is the same generate-and-verify architecture that appears in the
jailbreak analysis and the biosecurity screening system, applied to
software verification: neural generation of candidate proofs, symbolic
verification of their correctness.

The practical consequence: formal verification, previously expensive
enough to be reserved for avionics and nuclear systems, becomes
tractable for ordinary software as AI reduces the human effort
required. The cost barrier that kept the industry at level 1
(test-and-hope) is being lowered by the same technology that's
raising the cost of staying there.
\end{proposition}

%% =================================================================
\section{Cross-references}
%% =================================================================

\begin{remark}[Connections to other formalizations]
\begin{itemize}
\item \textbf{Un-Jailbreakable AI} (2026-07-10): generate-and-verify.
  AI-assisted formal verification uses the same architecture:
  neural generation of candidates, symbolic verification of
  correctness.
\item \textbf{Avoiding AGI Catastrophe Part 2} (2026-06-10): MPC and
  correct-by-construction. The ASI:Chain infrastructure described
  here is the same infrastructure that the MPC-based cryptographic
  laterality runs on.
\item \textbf{Decentralized Digital} (2026-08-06): ASI:Chain
  infrastructure. The correct-by-construction consensus mechanism
  is part of the broader ASI:Chain layer-1 development.
\item \textbf{Seeding RSI} (2026-07-28): MeTTa as implementation
  language. MeTTa's mathematical foundations are what make the
  correct-by-construction derivation tractable.
\item \textbf{How Decentralized AI Can Help Avert Bioterrorism}
  (2026-06-11): neural-symbolic screening. The AI-assisted
  verification described here uses the same neural-symbolic
  architecture applied to DNA sequence screening there.
\item \textbf{The Anthropic Fable Farce} (2026-07-08): capability
  diffusion. AI-accelerated bug discovery is capability diffusion
  applied to the vulnerability landscape --- the same base-space
  equalization dynamic.
\item \textbf{Non-invertible 2-cells} (LTM 85c02e32): cocycle defects.
  The Orchard bug is a concrete instance of the non-invertible
  witness infecting higher levels --- the local module correctness
  (1-cell) doesn't guarantee global composition correctness (2-cell).
\end{itemize}
\end{remark}

\begin{conjecture}[Verification reckoning timeline]
Conjecture: the software-verification reckoning will follow the
base-space-equalization timeline observed in AI capability
diffusion --- months, not years. The convergence will be driven
by AI-accelerated bug discovery making the cost of unverified
code exceed the cost of verification:
\[
  t_{\text{reckoning}} \approx t_0 + \Delta t \cdot
  \log\left(\frac{C_{\text{verify}}}{C_{\text{exploit}}}\right)
\]
where $C_{\text{verify}}$ is the cost of formal verification (decreasing
as AI improves) and $C_{\text{exploit}}$ is the expected cost of
exploitation (increasing as AI makes discovery cheaper). The reckoning
arrives when $C_{\text{verify}} < C_{\text{exploit}}$, at which point
the economic incentive flips from ``ship fast'' to ``ship proven.''
\end{conjecture}

\end{document}
